\chapter{微积分基本定理}
\label{chap:fundamental-theorem-calculus-solutions}

\begin{solution}{ftc-indefinite-step-function}
对
\[
 f(x)=0\ (x<0),\qquad f(x)=1\ (x\ge0)
\]
于 \([-1,1]\)，积分上限函数为 \(F(x)=0\)（\(x\le0\)）和
\(F(x)=x\)（\(x\ge0\)）。在 \(x<0\) 有 \(F'=0=f\)，在 \(x>0\)
有 \(F'=1=f\)。在 \(0\) 处左导数为 \(0\)、右导数为 \(1\)，所以不可导。
这恰与 \(f\) 只在 \(0\) 不连续相对应。
\end{solution}

\begin{solution}{ftc-continuity-point-general}
若 \(h>0\)，则
\[
 \abs{\frac{F(x_0+h)-F(x_0)}h-f(x_0)}
 \le\sup_{x_0\le t\le x_0+h}\abs{f(t)-f(x_0)}.
\]
若 \(h<0\)，反向积分给出
\[
 \frac{F(x_0+h)-F(x_0)}h
 =\frac1{\abs h}\int_{x_0+h}^{x_0}f(t)\,dt,
\]
故得到同样的上界，只需把区间改为
\([x_0+h,x_0]\)。连续性使两侧上界都趋于零。在 \(a\) 处只取
\(h>0\)，在 \(b\) 处只取 \(h<0\)，分别得到右导数和左导数。
\end{solution}

\begin{solution}{ftc-antiderivative-difference}
在任意闭子区间 \([x,y]\) 上，Lagrange 中值定理给出
\[
 (F-G)(y)-(F-G)(x)=(F-G)'(\xi)(y-x)=0.
\]
所以在一个区间上差为常数。若定义域是实线上的开集，每个连通分支都是
开区间，上述论证在每一分支上分别成立，不同分支的常数可以不同。
例如在 \((-2,-1)\cup(1,2)\) 上，零函数和分别取值0、1的函数有相同导数，
但其差不在整个定义域上恒定。
\end{solution}

\begin{solution}{ftc-integral-zero-continuous}
令 \(F(x)=\int_a^x f\)。题设说 \(F\equiv0\)。由微积分基本定理第一部分，
\(f=F'\equiv0\) 于开区间；端点由连续性得到同样结论。
\end{solution}

\begin{solution}{ftc-derivative-integrable-newton}
任取分割 \(P\)。对每段应用中值定理，选
\(\xi_i\in(x_{i-1},x_i)\) 使
\[
 F(x_i)-F(x_{i-1})=F'(\xi_i)\Delta x_i.
\]
相加得 \(F(b)-F(a)=S(F',\dot P)\)。取一列网格趋于零的分割，
Riemann 可积性使右端趋于 \(\int_a^bF'\)，即得公式。
\end{solution}

\begin{solution}{ftc-integrable-no-primitive-darboux}
可取
\[
 f(x)=\begin{cases}0,&x<0,\\1,&x\ge0,\end{cases}
\]
于 \([-1,1]\)。它单调，故 Riemann 可积。若 \(f=F'\)，Darboux 定理
要求 \(F'\) 在连接任意两点的区间上取遍中间值；但 \(f\) 不取
\((0,1)\) 中任何数，矛盾。
\end{solution}

\begin{solution}{integral-mean-value-weight-sign}
在 \([-1,1]\) 上取 \(f(x)=x\)、\(g(x)=x\)。则
\[
 \int_{-1}^1g=0,\qquad
 \int_{-1}^1fg=\int_{-1}^1x^2\,dx=\frac23.
\]
若中值公式成立，右端应为 \(f(\xi)\int g=0\)，矛盾。这里权函数变号。
\end{solution}

\begin{solution}{integral-average-between-extrema}
平均值总在 \([m,M]\) 中。若平均值等于 \(M\)，则连续非负函数
\(M-f\) 的积分为零，故 \(f\equiv M\)。等于 \(m\) 同理。因此非常值时
平均值严格介于两端。常函数显然在端点取平均值。
\end{solution}

\begin{solution}{hermite-hadamard}
记 \(m=(a+b)/2\)。对 \(x\in[a,m]\)，凸性给出
\[
 f(m)\le\frac{f(x)+f(a+b-x)}2.
\]
积分并在第二项中作代换 \(u=a+b-x\)，可得
\[
 (b-a)f(m)\le\int_a^b f(x)\,dx.
\]
又因
\[
 x=\frac{b-x}{b-a}a+\frac{x-a}{b-a}b,
\]
弦下估计给出
\[
 f(x)\le\frac{b-x}{b-a}f(a)+\frac{x-a}{b-a}f(b).
\]
在 \([a,b]\) 上积分，右端两个权函数的积分都等于 \((b-a)/2\)，故
\[
 \frac1{b-a}\int_a^bf(x)\,dx\le\frac{f(a)+f(b)}2.
\]
若 \(f\) 严格凸，则第一组点对在 \(x\ne m\) 时互异，第二个弦下估计在
\(a<x<b\) 时严格。两处差函数均连续、非负且在某个内点为正，
所以各自的积分为正，两个不等号都严格。
\end{solution}

\begin{solution}{integration-by-parts-log}
在连接 \(1,b\) 的闭区间上取 \(u=\log x\)、\(v'=1\)，于是 \(v=x\)。
两函数均为 \(C^1\)，故分部积分合法：
\[
 \int_1^b\log x\,dx
 =[x\log x]_1^b-\int_1^b1\,dx
 =b\log b-b+1.
\]
\end{solution}

\begin{solution}{substitution-nonmonotone}
若 \(f=1\)，左端为
\[
 \int_{-1}^1 2t\,dt=0,
\]
右端为 \(\int_{\phi(-1)}^{\phi(1)}1\,dx=\int_1^1 1\,dx=0\)。
若 \(f(x)=x\)，左端为
\[
 \int_{-1}^1 t^2(2t)\,dt=0,
\]
右端仍是从 \(1\) 到 \(1\) 的积分。参数从 \(1\) 降至 \(0\) 再升至
\(1\)，两段方向相反，贡献抵消。
\end{solution}

\begin{solution}{substitution-oriented-integral}
令 \(x=-u\)，即 \(\phi(u)=-u\)，并令 \(u\) 从 \(-a\) 到 \(-b\)。
代换公式给出
\[
 \int_{-a}^{-b}f(-u)(-1)\,du=\int_a^b f(x)\,dx,
\]
换元产生的负号与交换积分上下限产生的负号抵消，因此
\[
 \int_a^b f(x)\,dx
 =-\int_{-a}^{-b}f(-u)\,du
 =\int_{-b}^{-a}f(-u)\,du.
\]
特别取 \(a=0,b=1,f=1\)，左右两边都为 \(1\)，可用来检查符号。
若 \(f\) 为奇函数，换元给
\(\int_{-c}^0f(x)\,dx=\int_0^cf(-u)\,du=-\int_0^cf(u)\,du\)，
两半区间积分互相抵消；若为偶函数，
\[
 \int_{-c}^c f=2\int_0^c f.
\]
\end{solution}

\begin{solution}{ftc-taylor-integral-remainder-bound}
从积分余项，
\[
 \abs{R_{n,a}f(x)}
 \le\frac{M}{n!}\int_{\min(a,x)}^{\max(a,x)}
       \abs{x-t}^n\,dt
 =\frac{M}{(n+1)!}\abs{x-a}^{n+1}.
\]
这个论证只把整个积分用统一上界控制，没有证明积分等于某个中间点函数值
乘以长度；后者需要积分中值定理和相应连续性。
\end{solution}

\begin{solution}{integral-jensen-square}
对凸函数 \(\Phi(u)=u^2\) 和常权 \(w=1\) 应用积分 Jensen 不等式即得。
等号条件也可直接由
\[
 \frac1{b-a}\int_a^b(f-\bar f)^2
 =\frac1{b-a}\int_a^b f^2-\bar f^{\,2}
\]
判断，其中 \(\bar f=(b-a)^{-1}\int f\)。连续非负函数
\((f-\bar f)^2\) 积分为零当且仅当它恒为零，即 \(f\) 为常函数。
\end{solution}

\begin{solution}{ftc-absolute-value-antiderivative}
函数
\[
 F(x)=\frac12x\abs x
\]
在 \(x\ne0\) 的导数为 \(\abs x\)，在 \(0\) 的差商为
\(\frac12\abs h\to0\)，故 \(F'=\abs x\)，且导数连续，所以
\(F\in C^1\)。符号函数在 \(0\) 有跳跃，不具 Darboux 性质，因而不能
成为任何导函数。
\end{solution}

\begin{solution}{ftc-integral-upper-function-convex}
由基本定理 \(F'=f\)。若 \(f\) 单调不减，则可导函数的导数判据说明
\(F\) 凸。若 \(f\) 严格递增，则 \(F'\) 在任何非退化子区间上都不恒定，
因而 \(F\) 严格凸。在单调不减的假设下，严格递增等价于不在任何非退化
子区间上为常数：若 \(x<y\) 而 \(f(x)=f(y)\)，单调性便迫使
\(f\) 在 \([x,y]\) 上恒定。
\end{solution}

\begin{solution}{ftc-second-derivative-integral-reconstruction}
对函数 \(f\) 应用本章 Taylor 积分余项的 \(n=1\) 情形，直接得到
\[
 f(x)=f(a)+f'(a)(x-a)+\int_a^x(x-t)f''(t)\,dt.
\]
这里没有使用二重积分换序；所需公式已经由一元分部积分证明。

为证明凸性，须把展开中心改为任意 \(y\in[a,b]\)。同一公式给出
\[
 f(x)-f(y)-f'(y)(x-y)=\int_y^x(x-t)f''(t)\,dt.
\]
若 \(x\ge y\)，右端非负；若 \(x<y\)，将其写成
\(\int_x^y(t-x)f''(t)\,dt\)，仍非负。故函数位于每一条切线上方。
具体地，取 \(z=(1-\lambda)x+\lambda y\)、\(0<\lambda<1\)，
把以 \(z\) 为中心、分别在 \(x,y\) 处的两个切线不等式乘以
\(1-\lambda,\lambda\) 后相加，线性项消去，得到
\(f(z)\le(1-\lambda)f(x)+\lambda f(y)\)。
\end{solution}

\begin{solution}{integral-jensen-weighted-equality}
若 \(f\) 常值，直接代入即取等。以下设它非常值，记
\(l=\min f<\max f=r\)、\(W=\int_a^b w>0\)，并令
\(m=W^{-1}\int_a^b fw\)。连续非负函数 \((f-l)w\) 与
\((r-f)w\) 都在某个内点为正，故积分均为正，得到 \(l<m<r\)。
因此 \(m\) 是 \(\Phi\) 定义区间的内点，一维次梯度公式保证可取
\(p\in\partial\Phi(m)\)。令
\[
 H(u)=\Phi(u)-\Phi(m)-p(u-m)\ge0.
\]
若某个 \(u\ne m\) 有 \(H(u)=0\)，严格凸性会使连接 \(u,m\) 的
内点处 \(H<0\)，与支撑线不等式矛盾。所以 \(H(u)>0\) 对一切
\(u\ne m\) 成立。\(H(f(x))w(x)\) 连续非负，且由 \(f\) 非常值在
某个内点为正，故
\[
 0<\int_a^b H(f)w=\int_a^b\Phi(f)w-W\Phi(m).
\]
这证明非常值情形必严格，也完成等号条件的两向证明。
\end{solution}

\begin{solution}{ftc-derivative-discontinuity-primitive}
对 \(x\ne0\)，
\[
 F'(x)=2x\sin(1/x^2)-\frac2x\cos(1/x^2),
\]
而 \(F'(0)=\lim_{h\to0}h\sin(1/h^2)=0\)。第二项使 \(F'\) 在
\(0\) 附近无界：取 \(x_n=(2\pi n)^{-1/2}\)（\(n\ge1\)），则
\(F'(x_n)=-2/x_n\to-\infty\)；负半轴可取 \(-x_n\)。故不连续，
但它按定义仍是 \(F\) 的导函数。Riemann 可积函数必须有界，所以
\(F'\) 在任何包含 \(0\) 的非退化闭区间上都不
Riemann 可积。
\end{solution}
